\section{Polynomial interpolation}\label{sec:interp}
%
The proof of Theorem~\ref{thm:alg,interp} proceeds by an influential method known as polynomial interpolation introduced by Barvinok, see e.g.~\cite{Barvinok16}. 
In our application of this method, for some real $w^*>0$ we show that there is a region $\cU\subset\CC$ containing the interval $[1,w^*]$ on which $Z(G;w)\ne0$ for any UG instance $G$ of maximum degree $\Delta$. 
The region $\cU$ is a \emph{zero-free region}, and it guarantees that a Taylor series for (a suitable modification of) $\log Z$ converges inside $\cU$.
We then approximate $Z$ by computing the coefficients of a truncation of the Taylor series.
We require that $\cU$ is independent of the size of the graph $G$ to obtain a good approximation. 
The analysis yielding the approximation from $\cU$ is rather standard, e.g.~\cite{Barvinok16,Barvinok18,PR17}, though we include it in Appendix~\ref{app:interp,details} for completeness. The main technical work is in the following theorem establishing the region $\cU$. 
We write $\cN(S, \eta)$ for an open set in $\CC$ containing the open ball of radius $\eta$ around every point in $S$.
%
\begin{theorem}\label{thm:zero-free}
Let $k\in\NN_{\geq 3}$ and $\Delta\in \NN_{\geq 3}$.
Then with $w^*= 1+(\log k -1)/\Delta$ there exists $\eta=\omega(\frac{1}{\Delta\log k})$ such that for any $w\in \cN\left(\left[1,w^*\right],\eta\right)$ and any \UG[k] instance $G$ of maximum degree at most $\Delta$, $Z(G;w)\neq 0$.
\end{theorem}
%
%
Our proof of Theorem~\ref{thm:zero-free} is inductive in the style of~\cite{BDPR18} which gives a zero-free region for the antiferromagnetic Potts model, though here we have a generalization of the \emph{ferromagnetic} Potts model rather than than the \emph{antiferromagnetic} Potts model that was studied in~\cite{BDPR18}. 
%
%
Consider an $n$-vertex \UG[k] instance with $G=(V,E,\pi)$ with maximum degree $\Delta$.
In order to prove our results, we will need to work more generally with the partition function with boundary conditions.
For $m>0$ and a list $W=w_1\ldots w_m$ of distinct vertices of $V$ and a list $L=\ell_1\ldots\ell_m$ of pre-assigned colours in $[k]$ for the vertices in $W$ the \emph{restricted partition function} $Z^{W}_{L}(G;w)$ is defined by
\[
Z^{W}_{L}(G;w):=\sum_{\substack{\{x_u\}_{u\in V}\in[k]^V\\\{x_u\}_{u\in V} \text{ respects }(W,L)}}\;\prod_{\substack{(u,v)\in E,\\ x_v=\pi_{uv}(x_u)}}w\,,
\]
where we say that a color assignment $\{x_u\}_{u\in V}$ respects $(W,L)$ if for all $i=1\ldots, m$ we have $x_{w_i}=\ell_i$.
As it does not vary in the steps of the proof, we will omit the parameter $w$ and write $Z^W_L(G)$ for $Z^W_L(G;w)$.
We call the vertices $w_1,\ldots,w_m$ \emph{fixed} and refer to the remaining vertices in $V$ as \emph{free} vertices. The length of $W$ (resp.\ $L$), written $|W|$ (resp.\ $|L|$) is the length of the list.
Given a list of distinct vertices $W' = w_1\ldots w_m$, and a vertex $u$ (distinct from $w_1, \ldots, w_m$) we write $W = W'u$ for the concatenated list $W = w_1\ldots w_m u$ and we use similar notation $L'\ell$ for concatenation of lists of colours.  We write $\deg(v)$ for the degree of a vertex $v$ and we write $G\setminus uv$ ($G-u$) for the graph obtained from $G$ by removing the edge $uv$ (by removing the vertex $u$). 
%
%
%
To prove Theorem~\ref{thm:zero-free} we consider the same statement for restricted partition functions and induct over the number of vertices whose color is not fixed by the boundary conditions.
With a strengthened induction hypothesis we can argue that unfixing the specified color of a single vertex cannot affect the value of the partition function too much and continue the induction. 
The main technical difficulties are to bound the change in angle and radius unfixing a vertex can induce in the value of the partition function (as a complex number). 
%
\begin{restatable}{lemma}{zerofreeinduction}
\label{lem:kcol}
Let $\Delta\in \mathbb{N}_{\geq 3}$ and let $k\in \mathbb{N}_{\geq 3}$.  
Let $c=\log k-1$ and $\alpha=\log k^{1/2}-1$.
Then there exists constants $0<\eps<\theta<\frac{\pi}{3\Delta}$ with $\eps,\theta=\omega(1/\Delta)$ and $\eta=\omega(1/(\Delta\log k))$ such that for any $w\in \mathcal{N}([1,1+c/\Delta],\eta)$ and any \UG[k] instance $G$ of maximum degree at most $\Delta$
the following hold.
\begin{enumerate}
 \item\label{itm:kcol,nonzero}
 For all lists $W$ of distinct vertices of $G$ and all lists of pre-assigned colours $L$ of length $|W|$, $Z_L^W(G) \ne 0$.
 \item\label{itm:kcol,leaf}
 For all lists $W=W'u$ of distinct vertices of $G$ such that $u$ is a leaf and any two lists $L'l$, $L'l'$ of length $|W|$, the following hold.
 \begin{enumerate}
 \item\label{itm:kcol,leaf,free}
 If the unique neighbour $v$ of $u$ is free,
 \begin{enumerate}
 \item\label{itm:kcol,leaf,free,angle}
  the angle between vectors $Z_{L'l}^{W'u}(G)$ and $Z_{L'l'}^{W'u}(G)$ is at most $\theta$, and
 \item\label{itm:kcol,leaf,free,length}
       \begin{itemeq}
          \frac{Z_{L'l}^{W'u}(G)}{Z_{L'l'}^{W'u}(G)}   \leq 1 + \frac{\alpha}{\Delta}\,.
       \end{itemeq}
 \end{enumerate}
 \item\label{itm:kcol,leaf,fixed}
 If the unique neighbour $v$ of $u$ is fixed,
 \begin{enumerate}
 \item\label{itm:kcol,leaf,fixed,angle}
 the angle between vectors $Z_{L'l}^{W'u}(G)$ and $Z_{L'l'}^{W'u}(G)$ is at most $\eps$, and
 \item\label{itm:kcol,leaf,fixed,length} 
       \begin{itemeq}
          \frac{|Z_{L'l}^{W'u}(G)|}{|Z_{L'l'}^{W'u}(G)|}   \leq 1 + \frac{c}{\Delta}\,.
       \end{itemeq}
 \end{enumerate}
 \end{enumerate}
 \item\label{itm:kcol,nonleaf}
 For all lists $W = W'u$ of distinct vertices of $G$, and for all lists of pre-assigned colours $L'$ of length $|W'|$, let $d$ be the number of free neighbours of $u$, and let $b=\Delta-d$. Then for any pair of colours $l,l'$,
 \begin{enumerate}
 \item\label{itm:kcol,nonleaf,angle} 
 the angle between vectors $Z_{L'l}^{W'u}(G)$ and $Z_{L'l'}^{W'u}(G)$ is at most $d \theta + b\eps$, and 
 \item\label{itm:kcol,nonleaf,length}
       \begin{itemeq}
          \frac{|Z_{L'l}^{W'u}(G)|}{|Z_{L'l'}^{W'u}(G)|}   \leq(1+\alpha/\Delta)^d(1+c/\Delta)^{\Delta-d}\,.
       \end{itemeq}
 \end{enumerate} 
\end{enumerate}
\end{restatable}
%
Note that Statement~\ref{itm:kcol,nonzero} with $W=L=\emptyset$ is the result we want for Theorem~\ref{thm:zero-free}, Statement~\ref{itm:kcol,leaf} shows that changing the fixed color of a leaf (degree 1) vertex $u$ affects the angle and length of the restricted partition function by a small amount (depending on whether the neighbour of $u$ is itself free or fixed), and Statement~\ref{itm:kcol,nonleaf} is a similar but weaker version for any vertex. 
%
%
